\subsection{Limitations and threats to validity}
\label{sec:threats_to_validity}

This section distinguishes the principal threats to validity and the main limitations of the study.
Doing so is important because adversarial machine-learning results can appear more general or more conclusive than the underlying evidence really supports.
The present project provides a controlled empirical study, not a universal theory of one-pixel attacks or explanation robustness.
Its conclusions are therefore strongest when read at the level of the exact benchmark setting that was evaluated.

\subsection{Construct validity}
Construct validity concerns whether the quantities measured in the report adequately capture the concepts they are intended to represent.
This is especially important in the present project because both ``explanation robustness'' and ``guidance quality'' are indirect constructs rather than directly observable ground truths.

For explanation robustness, the report uses IoU@10, Spearman rank correlation, and deletion AUC.
These are all reasonable and informative metrics, but each captures only one aspect of a larger idea.
IoU@10 measures overlap of highly ranked pixels, not semantic equivalence of the full explanation.
Spearman correlation measures agreement in attribution ordering, not whether the explanation is human-meaningful.
Deletion AUC measures one operational notion of faithfulness, but it is still a proxy rather than a definitive measure of causal importance.
As a result, the report should be read as analysing measurable properties of explanation outputs under attack, not as proving exactly how the classifier internally reasons.

A second construct-validity issue concerns the use of RISE saliency as a responsibility map for guidance.
A RISE map is an estimate of which pixels are influential for the model's \emph{current} prediction.
That is not the same thing as a map of which pixel change yields the most query-efficient route to \emph{misclassification}.
Therefore, a negative result for RISE-guided DE should be interpreted narrowly: it shows that this specific black-box saliency prior does not improve the success--cost frontier in the tested configuration.
It does not prove that all explanation-guided one-pixel attacks must fail.

\subsection{Internal validity}
Internal validity concerns whether the reported outcomes genuinely arise from the experimental factors under study rather than from confounds in implementation, preprocessing, or aggregation.
Several design choices strengthen internal validity in this project.
The evaluated image range is fixed.
The eligibility rule is fixed.
The attack variants are compared on the same pretrained model and the same eligible subset.
The explanation metrics are computed through a stable pipeline, and the final report imports generated artefacts rather than copying numbers manually.

Even so, some internal-validity risk remains.
The project depends on a multi-stage pipeline involving attacks, saved tensors, explanation generation, summary scripts, and report imports.
Any complex experimental pipeline can contain silent implementation mistakes, schema mismatches, or device-specific quirks.
The report mitigates these risks through sanity checks, smoke tests, fixed seeds, and frozen artefact traceability rather than through formal verification.
That is appropriate for an individual research project, but it still means that the results ultimately rely on a carefully tested software artefact rather than a mathematically verified implementation.

Another internal-validity consideration is comparison symmetry.
Fastprior uses a fixed-budget multi-target design, while the untargeted and RISE-guided variants use adaptive stopping.
This makes the comparison analytically useful, but not perfectly matched in optimisation structure.
The report partially addresses this by separating median cost over all eligible images from median cost over successful images and by interpreting fastprior as a deliberately different trade-off point rather than as a drop-in replacement for untargeted DE.

\subsection{External validity}
The strongest limitation of the study is external validity.
All main results are derived from a single dataset, CIFAR-10, and a single main classifier architecture, a pretrained ResNet-20.
The perturbation budget is restricted to one pixel, and the threat model is black-box.
These design choices are defensible because they produce a clear and tractable benchmark, but they also limit generalisation.

The findings may not transfer directly to larger images, different architectures, different datasets, or other perturbation budgets.
For example, a one-pixel perturbation on a $32\times 32$ CIFAR-10 image is a very different object from a one-pixel perturbation on a higher-resolution natural image.
Likewise, responsibility-map guidance might behave differently in settings where spatial structure is richer, where the model architecture changes the saliency landscape, or where more than one pixel can be modified.
The report therefore supports claims about a controlled benchmark setting, not about universal behaviour across computer vision systems.

\subsection{Selection and composition effects}
A particularly important limitation of the present design is that explanation metrics are computed only on successful adversarial examples.
This is appropriate for the research questions, especially RQ1, but it introduces a selection effect.
The explanation statistics do not describe all attacked images.
They describe the subset on which the attack succeeds.

This matters because different attack variants can succeed on different kinds of images.
A guided attack may appear to produce more stable explanations not because the perturbation mechanism is inherently better, but because it succeeds on a narrower family of examples whose explanation maps are already less volatile under sparse perturbation.
The report explicitly acknowledges this interpretation in the discussion chapter and avoids treating explanation stability improvements as automatically equivalent to a better attack.
Still, the selection effect remains a genuine limitation, and it is one reason the explanation comparisons must be interpreted with caution.

\subsection{Measurement and aggregation limitations}
The report uses medians for attack cost because query distributions are skewed.
This is a sensible choice, but medians do not fully describe distribution shape.
Two runs can share the same median while differing in tail behaviour, variance, or failure-mode structure.
Similarly, mean explanation metrics provide compact summaries, but they can hide heterogeneity across images or classes.

The report partly compensates for this by reporting both means and medians for several key explanation quantities and by separating all-image versus successful-image cost summaries.
However, the document still does not provide a full distributional analysis of costs or explanation metrics.
That choice keeps the report focused, but it limits how much can be said about rare or extreme cases.

\subsection{Summary of limitations}
The main limitations of the study can therefore be summarised as follows:
\begin{itemize}
  \item explanation metrics are informative proxies rather than direct measurements of internal model reasoning;
  \item the guidance result applies to one specific saliency-derived prior rather than to all possible guidance mechanisms;
  \item the evaluation is limited to one dataset, one main architecture, and one sparse black-box threat model;
  \item successful-only explanation analysis introduces genuine selection and composition effects;
  \item summary statistics do not exhaustively characterise the full distributions of attack cost and explanation behaviour.
\end{itemize}

These limitations do not invalidate the project.
Rather, they define the correct level of interpretation.
The report provides a controlled empirical account of one-pixel attacks, explanation robustness, and guidance trade-offs in a reproducible benchmark setting.
Its conclusions are strongest within that scope.
